% ============================================================
% PHẦN B — KIẾN THỨC NỀN SCIMKG — SƯỜN
% Định nghĩa độc lập để đọc Phần B không cần quay lại Phần A.
% ============================================================
\section{SciMKG: kiến thức nền}\label{sec:sm-background}

\subsection{Đồ thị tri thức giáo dục và đa phương thức}
\todo{Định nghĩa đồ thị tri thức bằng các bộ ba (đầu, quan hệ, đuôi); đồ thị tri thức giáo dục (EKG), điểm kiến thức và khái niệm; đồ thị tri thức đa phương thức gắn khái niệm với văn bản, ảnh, video và âm thanh. ConceptNet là nguồn bổ sung khái niệm trong SciMKG.}

\subsection{LLM cho trích xuất thông tin}
\todo{Trích xuất khái niệm và nhận dạng thực thể (NER) bằng LLM; GPT-NER, UniversalNER; câu lệnh đầu vào quy định miền kiến thức và định dạng đầu ra. Phân biệt gọi mô hình có sẵn lúc suy luận với tinh chỉnh trọng số.}

\subsection{Tự phản hồi và bỏ phiếu nhiều lần sinh}
\todo{Self-Refine: vòng sinh phản hồi rồi sửa đầu ra; Self-Consistency: tổng hợp nhiều lần sinh bằng bỏ phiếu. Giải thích vai trò của nhiều LLM, số vòng và ngưỡng giữ khái niệm trong EVIA.}

\subsection{Các chỉ số đánh giá trích xuất}
\todo{Định nghĩa Precision, Recall, F1, độ chính xác (ACC), MCC và AUC; nêu rõ đơn vị đánh giá, nhãn vàng, tập mẫu âm và điểm số cần để tính AUC.}

\subsection{Căn chỉnh ảnh--văn bản}
\todo{Bài toán gắn ảnh với khái niệm; điểm căn chỉnh ảnh--văn bản CLIPScore và BLIP-ITM; phân biệt điểm tự động với đánh giá nội dung của chuyên gia.}
